\chapter{Conclusions}\label{chap:conclusions}

% ============================================================================
\section{Summary of Contributions}\label{sec:summary}
% ============================================================================

This thesis set out to close a specific gap in cybersecurity education: the absence of an accessible, browser-based environment that simulates the complete workflow of hardware analysis, from physical board inspection to embedded system exploitation. We proposed and implemented \pcbctf{}, a web-based interactive cyber range for hardware security training in Capture The Flag format. The design was guided by two overarching principles---\emph{interaction fidelity} and \emph{declarative configurability}---and yielded four contributions that we now revisit against the goals stated in Chapter~\ref{chap:intro}.

The first contribution, the client-side instrument simulation architecture, delivers realistic hardware analysis tools---digital multimeter, UART adapter, SPI programmer---directly in the browser. Differential voltage measurements, crossover wiring validation, and role-based probe connections reproduce the diagnostic reasoning a student would exercise on a physical workstation. A normalised coordinate system ensures that interactive overlays scale correctly across display sizes, and the absence of server-side execution eliminates an entire class of security concerns while enabling offline operation after initial page load.

The second contribution, the declarative terminal emulation system, provides a configuration language expressive enough to define virtual filesystems, multi-stage boot sequences, custom commands with six output types, five constraint categories, conditional flag unlock side effects, and multi-tab environments with independent boot sequences. This system reproduces the experience of interacting with a real embedded device through a synchronous command pipeline. Exercise authors shape terminal behaviour entirely through JSON configuration, never through code changes.

The third contribution, the no-code authoring interface, provides a visual exercise editor with drag-and-drop PCB markup, structured step and objective editing, integrated terminal configuration, and a preset system for exercise versioning and distribution. The four-store state management architecture (\texttt{exerciseStore}, \texttt{settingsStore}, \texttt{terminalSettingsStore}, \texttt{presetStore}) supports automatic dirty tracking, three-level persistence, and live synchronisation between the authoring panel and the simulator.

The fourth contribution, the validation case study, demonstrates that a complete multi-step exercise---covering component identification, UART connection, firmware extraction, and six distinct vulnerability discovery tasks on a TP-Link WR841N router---can be fully configured and delivered through the proposed system.

% ============================================================================
\section{Requirements Assessment}\label{sec:req-assessment}
% ============================================================================

The functional and non-functional requirements defined in Chapter~\ref{chap:contributions} provide a concrete basis for evaluating the platform. All twenty-three functional requirements (RF-01 through RF-23) are satisfied by the current implementation: the PCB simulation subsystem (RF-01 through RF-03) is validated by the interactive viewer described in Chapter~\ref{chap:implementation}; the measurement and connection tools (RF-04 through RF-07) are exercised in the case study's multimeter, UART, and SPI tasks; the terminal subsystem (RF-08 through RF-13) supports the full range of commands, constraints, and boot sequences required by the WR841N exercise; the exercise structure (RF-14 through RF-17) governs the ten-step learning path; and the authoring subsystem (RF-18 through RF-23) enables the complete exercise to be created, modified, and distributed without code changes. Among the non-functional requirements, zero installation (RNF-01), client-side execution (RNF-06), server independence (RNF-07), and the prohibition of server-side code execution (RNF-10) are architectural invariants enforced by the implementation's design. The remaining non-functional requirements---simplified deployment (RNF-02), guided interaction (RNF-03), visual feedback (RNF-04), responsive layout (RNF-05), static typing (RNF-08), modular terminal design (RNF-09), and session isolation (RNF-11)---are met as documented in Chapters~\ref{chap:architecture} and~\ref{chap:implementation}.

% ============================================================================
\section{Discussion}\label{sec:discussion}
% ============================================================================

The comparative analysis in Chapter~\ref{chap:stateoftheart} positions \pcbctf{} in a space that existing solutions do not occupy: a zero-cost, fully browser-based platform offering PCB-level interaction, instrument simulation, and terminal emulation with no server-side code execution. This combination addresses the three main barriers to hardware security training---cost (no equipment required), logistics (no physical setup or shipping), and security (no risk of server compromise through student inputs).

The pedagogical design of the platform draws on well-established principles. Svabensky et al.'s survey of hands-on cybersecurity training confirms that practical, interactive exercises produce measurably better learning outcomes than passive instruction~\cite{svabensky2020}. The progressive structure of \pcbctf{} exercises---where each step builds on skills acquired in previous steps, new tools become available only when the student has demonstrated competence with existing ones, and hints are available but not forced---embodies the concept of \emph{scaffolding}, in which instructional support is calibrated to the learner's current level and gradually withdrawn as competence grows. Similarly, the progressive flag system implements \emph{progressive disclosure}: flag fragments are revealed only when the student demonstrates understanding of the underlying concept (recognising a boot configuration anomaly, cracking a password hash, identifying a command injection pattern), reinforcing the connection between analytical reasoning and tangible outcomes. The student cannot guess the flag; each fragment requires completing a specific analytical step, and the partial flag string provides continuous feedback on progress. Hamari et al.\ established that gamification elements---points, progressive unlocks, visible achievement markers---significantly increase engagement in learning contexts~\cite{hamari2014}, and the CTF format exploits precisely these mechanisms.

The pedagogical coverage of the platform spans the core phases of hardware analysis: visual reconnaissance, electrical measurement, protocol identification, serial connection, firmware extraction, bootloader analysis, credential extraction, and vulnerability discovery. This progression aligns with the methodology described in standard references such as the \emph{Hardware Hacking Handbook}~\cite{jasper2020hardware} and \emph{Practical IoT Hacking}~\cite{gupta2021practical}. The ten-step learning path of the case study exercise maps directly to the procedural phases described in Section~\ref{subsec:methodology}, demonstrating that the platform supports a structurally complete training experience within a single exercise.

The choice to simulate instrument readings with realistic electrical characteristics---voltage levels, impedances, noise---rather than presenting idealised values reflects a commitment to developing transferable diagnostic reasoning. When a student measures 3.3\,V with low impedance on a UART TX line and 3.3\,V with high impedance on the RX line, she is practising the same inference she would apply on real hardware. The simulated digital noise on the multimeter display, while cosmetic, reinforces the perception of working with a physical instrument.

The deterministic nature of the simulation is both a strength and a limitation. It is a strength because identical conditions across all sessions make the platform suitable for graded assessments and reproducible demonstrations. It is a limitation because real hardware analysis involves noise, timing variations, and unexpected device behaviour that a configuration-driven simulation cannot fully capture. We therefore position the platform as a preparatory tool that develops foundational reasoning skills, not as a replacement for physical laboratory experience.

The security-by-design approach documented in Section~\ref{sec:security} demonstrates that the platform practises what it teaches: by confining all execution to the client side, sanitising inputs, and enforcing path boundaries, the implementation avoids the very classes of vulnerabilities---injection, path traversal, insecure defaults---that its exercises train students to identify in IoT devices.

% ============================================================================
\section{Limitations}\label{sec:limitations}
% ============================================================================

Five limitations bound the current implementation, each representing a conscious design trade-off rather than an oversight.

First, the terminal supports only pre-configured commands. A student who types a command absent from the exercise configuration receives a ``command not found'' error, even when the command would be meaningful in a real environment---\texttt{ifconfig}, \texttt{dmesg}, or \texttt{mount}, for example. This restricts the sense of open-ended exploration that characterises interaction with a real device. The trade-off is deliberate: supporting only declared commands guarantees deterministic behaviour and eliminates the risk of students encountering undefined states, both of which are essential for graded assessments. In practice, the limitation can be mitigated by the exercise author, who can pre-configure any command that is pedagogically relevant. The case study exercise includes over twenty custom commands, which proved sufficient to sustain a convincing exploration experience across six challenges.

Second, protocol coverage is limited to UART and SPI. JTAG, I\textsuperscript{2}C, and CAN---all common in embedded devices---are not yet supported. Each new protocol would require a dedicated tool widget, a pin type definition with role-based validation, and a set of corresponding exercise objective types. The current architecture is designed to accommodate such extensions: the modular tool system, the generic pin model, and the objective type registry are all extensible without structural changes. We prioritised UART and SPI because they cover the two most prevalent attack vectors in consumer IoT devices (serial console access and firmware extraction, respectively), as documented in Section~\ref{subsec:exploitability}.

Third, the platform operates in a single-user architecture without scoring, leaderboards, or persistent progress tracking. Each browser session is independent, and there is no concept of user identity or session history. This limits the platform's applicability to competitive CTF events, where simultaneous participation and ranking are expected. The single-user model was chosen to minimise architectural complexity and avoid the security implications of user authentication in an educational tool. Multi-user support is the most frequently cited extension in our future work plan.

Fourth, no empirical user evaluation has been conducted. The validation presented in Chapter~\ref{chap:casestudy} is based on feature coverage, structural alignment with real-world methodology, and informal observation during development. A controlled study with a representative student cohort---measuring learning outcomes through pre- and post-assessments, comparing against traditional laboratory instruction, and evaluating usability through standardised questionnaires---remains necessary to quantify pedagogical effectiveness. We acknowledge this as the most important outstanding task.

Fifth, the authoring panel is protected only by a build-time environment flag and lacks authentication or authorisation layers. While this is sufficient for the intended deployment model (single-user, local development), it makes the platform unsuitable for multi-tenant production environments where multiple instructors might share a single instance. Adding role-based access control to the authoring API is a straightforward extension that we defer to a future version.

% ============================================================================
\section{Future Work}\label{sec:future}
% ============================================================================

Several directions for future development follow naturally from these limitations and from the broader design goals of the platform.

\textbf{Multi-user support and scoring.} Adding user authentication, session management, database-backed progress persistence, and a timing-based leaderboard would enable the platform to host live CTF events with concurrent participants. A scoring system could incorporate partial credit for discovered flag fragments, time penalties for hint usage, and bonus points for optional objectives. This extension would also enable instructors to collect anonymised aggregate data on exercise difficulty and common stumbling points.

\textbf{Hybrid simulation--hardware integration.} A natural evolution connects the platform to remote laboratory infrastructure, enabling a seamless transition from simulated exercises to interaction with real hardware. A student might begin with the simulated WR841N exercise to develop foundational skills, then switch to a remote lab session with web-controlled instruments targeting a physical device. The platform's clean separation between simulation frontend and configuration backend is amenable to this extension: instrument widgets could relay commands to remote hardware rather than to local state.

\textbf{Additional protocol support.} Implementing JTAG would enable exercises focused on boundary scan, debug port exploitation, and memory read-out. I\textsuperscript{2}C support would extend coverage to EEPROM-based devices. CAN bus support would open the door to automotive security scenarios, a rapidly growing area of concern. Each new protocol requires a pin type definition, a tool widget, and a validation strategy in the exercise store.

\textbf{Dynamic terminal responses.} The restriction to pre-configured commands could be relaxed by generating plausible output for unexpected commands based on the current context---filesystem state, boot stage, previous command history. Such a mechanism would improve the sense of open-ended exploration while maintaining the security guarantee of client-side-only execution.

\textbf{Student analytics and adaptive difficulty.} Tracking completion times, hint usage frequency, error patterns, and command sequences would provide instructors with quantitative data to calibrate exercise difficulty. An adaptive system could dynamically adjust hint visibility, present additional guidance for struggling students, or unlock bonus challenges for advanced learners.

\textbf{Empirical validation.} The most important outstanding task is a rigorous empirical evaluation. A controlled study with a representative student cohort---ideally within a university course---would quantify pedagogical effectiveness through pre- and post-assessments, measure usability through standardised questionnaires such as the System Usability Scale, and compare learning outcomes against traditional laboratory instruction. Such a study would provide the evidence base needed to justify broader adoption.

\bigskip

Hardware security knowledge should not be gated by the price of a logic analyser or the availability of a supervised laboratory. The barriers that make hardware skills scarce are primarily logistical, not intellectual, and a well-designed simulation platform can lower them dramatically. \pcbctf{} demonstrates that the core workflow of hardware analysis---from visual board inspection through protocol-level probing to embedded system exploitation---can be faithfully reproduced in a standard web browser at zero marginal cost per student. If the cybersecurity community is serious about closing the hardware skills gap, tools like this one are not optional; they are the necessary bridge between the software-centric training ecosystem that exists today and the hardware-aware workforce that tomorrow's threat landscape demands.
